% Local editorial examples following CP-VI-0075; canonical problem and solution remain unchanged.
\par\medskip
\noindent\textbf{Further explanation: the tensor product is the universal ring.}
An $A$-algebra homomorphism $B\otimes_A C\to R$ is equivalent to a
pair of $A$-algebra homomorphisms $B\to R$ and $C\to R$ whose images
commute.  Since $R$ is commutative, the commuting condition is automatic.
Applying this with $R=\Gamma(T,\mathcal O_T)$ yields naturally in $T$
\[
\operatorname{Hom}\!\bigl(T,\operatorname{Spec}(B\otimes_A C)\bigr)
\cong
\operatorname{Hom}(T,\operatorname{Spec} B)
\mathbin{\times}_{\operatorname{Hom}(T,\operatorname{Spec} A)}
\operatorname{Hom}(T,\operatorname{Spec} C).
\]
This also proves the asserted isomorphism for non-affine test schemes $T$.

\par\medskip
\noindent\textbf{Worked example 1 (a non-reduced intersection).}
Let $k$ be a field, $A=k[t]$, $B=k[x]$ with $t\mapsto x^2$,
and $C=k$ with $t\mapsto0$.  The relation identifying the two
images of $t$ gives
\[
B\otimes_A C
\cong k[x]\otimes_{k[t]} k
\cong k[x]/(x^2).
\]
The fiber product has one underlying point, but the class of $x$
is a nonzero nilpotent; it cannot be replaced by $\operatorname{Spec} k$.

\par\medskip
\noindent\textbf{Worked example 2 (two distinct points after base change).}
In the same example, send $t\mapsto a\in k^\times$ instead of $0$.
Then
\[
B\otimes_A k\cong k[x]/(x^2-a).
\]
If $\operatorname{char}k\ne2$ and $a=b^2$ for some $b\in k^\times$,
the factors $(x-b)$ and $(x+b)$ are comaximal, so the Chinese
remainder theorem yields $k[x]/(x^2-a)\cong k\times k$.
Consequently the fiber product consists of two reduced $k$-points.

\par\medskip
\noindent\textbf{Worked example 3 (two infinitesimal directions).}
Put $A=k$, $B=k[\varepsilon]/(\varepsilon^2)$, and
$C=k[\delta]/(\delta^2)$.  Then
\[
B\otimes_k C
\cong k[\varepsilon,\delta]/(\varepsilon^2,\delta^2),
\]
with $k$-basis $1,\varepsilon,\delta,\varepsilon\delta$.
The mixed infinitesimal $\varepsilon\delta$ survives: it is not
killed by either square-zero relation individually.

\par\medskip
\noindent\textbf{Worked example 4 (reduction modulo a prime).}
For $A=\mathbb Z$, $B=\mathbb Z[x]$, and $C=\mathbb F_p$,
\[
B\otimes_{\mathbb Z}\mathbb F_p
\cong\mathbb F_p[x].
\]
Thus changing the base from $\operatorname{Spec}\mathbb Z$ to
$\operatorname{Spec}\mathbb F_p$ changes the coefficient ring,
not the number of polynomial variables.  This is a basic instance
of extension of scalars.
